L'any 2004 es va aconseguir mesurar la massa d'un virus. Es va determinar la freqüència d'oscil·lació d'un braç horitzontal petitíssim, primer sense el virus i després amb el virus adherit. Sense el virus, la freqüència d'oscil·lació era de $2,00\times 10^{15}\ \mathrm{Hz}$, i amb el virus, aquesta freqüència va ser de $2,87\times 10^{14}\ \mathrm{Hz}$.

\begin{apartat}{1}
Si suposem que el braç horitzontal sense el virus adherit es comporta com una molla amb una massa oscil·lant de $2,10\times 10^{-16}\ \mathrm{g}$ lligada a un extrem, quina és la constant elàstica d'aquesta suposada molla?
\end{apartat}

\begin{apartat}{1}
Partint de la mateixa suposició anterior sobre el comportament oscil·latori del sistema, calculeu la massa del virus.
\end{apartat}
